\documentclass[11pt,a4paper]{article}

\usepackage[T1]{fontenc}
\usepackage[utf8]{inputenc}
\usepackage[french]{babel}
\usepackage{amsmath,amssymb,amsthm}

\newtheorem{lemma}{Lemme}

\newcommand{\Res}{\operatorname{Res}}

\begin{document}

\subsection{Preuve de la loi de réciprocité quadratique à l'aide des résultants}

Soient \(p\) et \(q\) deux nombres premiers impairs distincts. Posons
\[
m=\frac{p-1}{2}
\qquad\text{et}\qquad
n=\frac{q-1}{2}.
\]
Nous allons construire des polynômes \(Q_p,Q_q\in\mathbb{Z}[X]\) de degrés
respectifs \(m\) et \(n\), puis utiliser la relation
\[
\Res(P,Q)
=
(-1)^{\deg(P)\deg(Q)}\Res(Q,P).
\]

\subsubsection{Premier lemme}

\begin{lemma}
Pour tout nombre premier impair \(p\), il existe un polynôme unitaire
\(Q_p\in\mathbb{Z}[T]\), de degré \((p-1)/2\), tel que
\[
1+X+\cdots+X^{p-1}
=
X^{(p-1)/2}Q_p\left(X+\frac{1}{X}\right).
\]
De plus,
\[
Q_p(T)\equiv (T-2)^{(p-1)/2}\pmod p
\]
et
\[
Q_p(2)=p.
\]
\end{lemma}

\begin{proof}
Posons \(m=(p-1)/2\). Après division par \(X^m\), l'identité recherchée
s'écrit
\[
Q_p\left(X+\frac{1}{X}\right)
=
X^m+X^{m-1}+\cdots+X+1
+X^{-1}+\cdots+X^{-m}.
\]

Considérons le polynôme symétrique
\[
S_m(U,V)
=
U^m+U^{m-1}+\cdots+U+1+V+\cdots+V^{m-1}+V^m.
\]
Comme \(S_m\) est symétrique en \(U\) et \(V\), il existe
\(R_m\in\mathbb{Z}[A,B]\) tel que
\[
S_m(U,V)=R_m(U+V,UV).
\]
En prenant \(U=X\) et \(V=X^{-1}\), on obtient
\[
S_m\left(X,X^{-1}\right)
=
R_m\left(X+\frac{1}{X},1\right).
\]
Il suffit donc de poser
\[
Q_p(T)=R_m(T,1).
\]
Le polynôme \(Q_p\) est unitaire de degré \(m\).

Montrons maintenant la congruence modulo \(p\). Dans
\(\mathbb{F}_p[X]\), on a
\[
1+X+\cdots+X^{p-1}
=
\frac{X^p-1}{X-1}
\equiv
\frac{(X-1)^p}{X-1}
=
(X-1)^{p-1}.
\]
Comme \(p-1=2m\), il vient
\begin{align*}
1+X+\cdots+X^{p-1}
&\equiv (X-1)^{2m} \\
&=(X^2-2X+1)^m \\
&=X^m\left(X+\frac{1}{X}-2\right)^m
\pmod p.
\end{align*}
D'autre part,
\[
1+X+\cdots+X^{p-1}
=
X^mQ_p\left(X+\frac{1}{X}\right).
\]
Puisque \(X\) est inversible dans
\(\mathbb{F}_p[X,X^{-1}]\), on en déduit
\[
Q_p\left(X+\frac{1}{X}\right)
\equiv
\left(X+\frac{1}{X}-2\right)^m
\pmod p.
\]

L'application
\[
\begin{aligned}
\mathbb{F}_p[T]&\longrightarrow\mathbb{F}_p[X,X^{-1}],\\
F(T)&\longmapsto F\left(X+\frac{1}{X}\right)
\end{aligned}
\]
étant injective, on obtient
\[
Q_p(T)\equiv(T-2)^m\pmod p.
\]

Enfin, en évaluant l'identité initiale en \(X=1\), on trouve
\[
p=Q_p(2).
\]
\end{proof}

\subsubsection{Deuxième lemme}

\begin{lemma}
Pour deux nombres premiers impairs distincts \(p\) et \(q\), on a
\[
\Res(Q_p,Q_q)=\pm1.
\]
\end{lemma}

\begin{proof}
Montrons d'abord que \(Q_p\) et \(Q_q\) n'ont aucune racine commune dans
\(\mathbb{C}\).

Supposons qu'il existe \(\beta\in\mathbb{C}\) telle que
\[
Q_p(\beta)=Q_q(\beta)=0.
\]
Choisissons \(x\in\mathbb{C}^{\times}\) vérifiant
\[
x+\frac{1}{x}=\beta,
\]
c'est-à-dire
\[
x^2-\beta x+1=0.
\]
D'après le premier lemme,
\[
1+x+\cdots+x^{p-1}
=
x^mQ_p(\beta)
=
0.
\]
En multipliant par \(x-1\), on obtient
\[
x^p=1.
\]
De même,
\[
x^q=1.
\]
Comme \(\gcd(p,q)=1\), il vient \(x=1\). Mais alors
\[
1+x+\cdots+x^{p-1}=p\neq0,
\]
ce qui est contradictoire. Ainsi,
\[
\Res(Q_p,Q_q)\neq0.
\]

Posons
\[
R=\Res(Q_p,Q_q)\in\mathbb{Z}.
\]
Supposons qu'un nombre premier \(\ell\) divise \(R\). Les réductions de
\(Q_p\) et \(Q_q\) modulo \(\ell\) possèdent alors une racine commune
\[
\beta\in\overline{\mathbb{F}}_{\ell}.
\]
Choisissons \(x\in\overline{\mathbb{F}}_{\ell}\) tel que
\[
x^2-\beta x+1=0.
\]
Puisque le terme constant de ce polynôme est \(1\), on a \(x\neq0\), et
\[
\beta=x+\frac{1}{x}.
\]

D'après les identités définissant \(Q_p\) et \(Q_q\),
\[
1+x+\cdots+x^{p-1}
=
x^m\overline{Q}_p(\beta)
=
0
\]
et
\[
1+x+\cdots+x^{q-1}
=
x^n\overline{Q}_q(\beta)
=
0,
\]
où \(\overline{Q}_p\) et \(\overline{Q}_q\) désignent les réductions
respectives de \(Q_p\) et \(Q_q\) modulo \(\ell\).

En multipliant respectivement par \(x-1\), on obtient
\[
x^p=1
\qquad\text{et}\qquad
x^q=1.
\]
Comme \(\gcd(p,q)=1\), on en déduit \(x=1\). Les deux sommes précédentes
donnent alors
\[
p=0
\qquad\text{et}\qquad
q=0
\]
dans \(\overline{\mathbb{F}}_{\ell}\). Ainsi,
\[
\ell\mid p
\qquad\text{et}\qquad
\ell\mid q,
\]
ce qui est impossible puisque \(p\neq q\).

Aucun nombre premier ne divise donc l'entier non nul \(R\). Par conséquent,
\[
R=\pm1.
\]
\end{proof}

\subsubsection{Conclusion}

Posons
\[
R=\Res(Q_p,Q_q).
\]
D'après le deuxième lemme,
\[
R\in\{-1,1\}.
\]

Réduisons d'abord le résultant modulo \(p\). Le premier lemme donne
\[
Q_p(T)\equiv(T-2)^m\pmod p.
\]
Comme le résultant est une expression polynomiale à coefficients entiers
en les coefficients des deux polynômes, on obtient
\begin{align*}
R
&\equiv
\Res\left((T-2)^m,Q_q(T)\right) \\
&=
Q_q(2)^m \\
&=
q^m
\pmod p.
\end{align*}
Par le critère d'Euler,
\[
q^m
=
q^{(p-1)/2}
\equiv
\left(\frac{q}{p}\right)
\pmod p.
\]
Ainsi,
\[
R\equiv\left(\frac{q}{p}\right)\pmod p.
\]
Les deux membres appartenant à \(\{-1,1\}\), on en déduit
\[
R=\left(\frac{q}{p}\right).
\]

En échangeant les rôles de \(p\) et de \(q\), on obtient de même
\[
\Res(Q_q,Q_p)
=
\left(\frac{p}{q}\right).
\]
Or les degrés de \(Q_p\) et \(Q_q\) sont respectivement \(m\) et \(n\), donc
\[
\Res(Q_q,Q_p)
=
(-1)^{mn}\Res(Q_p,Q_q).
\]
Par conséquent,
\[
\left(\frac{p}{q}\right)
=
(-1)^{mn}\left(\frac{q}{p}\right).
\]
Puisque
\[
m=\frac{p-1}{2}
\qquad\text{et}\qquad
n=\frac{q-1}{2},
\]
on obtient finalement
\[
\boxed{
\left(\frac{p}{q}\right)
\left(\frac{q}{p}\right)
=
(-1)^{\frac{p-1}{2}\frac{q-1}{2}}
}.
\]
Ce qui correspond à la loi de la réciprocité quadratique.
\end{document}
